\chapter{Indeks pól wektorowych i twierdzenie Poincarégo--Hopfa z brzegiem}
\label{ch:poincare-hopf-v11}

\section{Dlaczego nie powtarzamy dowodu dla rozmaitości zamkniętych?}
Twierdzenie Poincarégo--Hopfa dla zamkniętych rozmaitości
zorientowanych, wraz z definicją lokalnego indeksu pola,
zostało już udowodnione w
Sekcji~\ref{sec:poincare-hopf-16}.
Z kolei Twierdzenie~\ref{thm:lefschetz-hopf-v10} utożsamia
liczbę Lefschetza z sumą lokalnych indeksów punktów stałych.
W tym rozdziale dokładnie porównamy oba rodzaje indeksów,
a potem potraktujemy rozmaitości z brzegiem.
Nie wolno bez komentarza stosować wzoru dla rozmaitości zamkniętych
do pola, które znika lub jest styczne do brzegu: pojawia się wtedy
dodatkowy wkład brzegowy.

\section{Od zera pola do punktu stałego}
\begin{lemat}[Lokalne porównanie indeksów]
\label{lem:field-fixed-v11}
Niech $M$ będzie gładką rozmaitością wymiaru $n\ge1$,
$X$ gładkim polem wektorowym, a $p\in M$ jego izolowanym zerem.
Wybierzmy metrykę riemannowską i otoczenie $U$ punktu $p$
o zwartym domknięciu, zawarte w jednej mapie, takie że
$X$ nie znika na $\partial U$. W razie potrzeby przedłużamy metrykę i pole do otoczenia
$M$ w rozmaitości bez brzegu, po czym dla dostatecznie
małego $t>0$ określamy lokalnie
$f_t(x)=\exp_x(-tX(x))$.
Definicja indeksu dotyczy wyłącznie wewnętrznego
zera $p$, więc nie zależy od wyboru przedłużenia.
Wówczas $f_t$ ma w $U$ jedyny punkt stały $p$, po odpowiednim
zmniejszeniu $U$, i
$\operatorname{ind}_p(f_t)=\operatorname{ind}_p(X)$.
Jeśli zero jest nieosobliwe, oba indeksy są równe
$\operatorname{sgn}\det(\D_pX)$.
\end{lemat}
\begin{proof}
Przyjmijmy mapę $\varphi$ z $\varphi(p)=0$.
Zwartość $\overline U$ pozwala wybrać $t>0$ tak małe,
aby dla każdego $x\in\overline U$ wektor $-tX(x)$ leżał
w obszarze, w którym $\exp_x$ jest różnowartościowe.
Stąd $f_t(x)=x$ wtedy i tylko wtedy, gdy $X(x)=0$.
Oznaczmy $F_t=\varphi\circ f_t\circ\varphi^{-1}$ i
$V=\mathrm d\varphi(X)$.
Na zwartym $\varphi(\overline U)$ zachodzi jednostajnie
$F_t(y)=y-tV(y)+O(t^2)$, toteż
$t^{-1}(y-F_t(y))\to V(y)$ jednostajnie.
Na $\varphi(\partial U)$ wektor $V$ jest oddalony od zera.
Wobec tego dla małego $t$ homotopia
$(1-s)V(y)+s\,t^{-1}(y-F_t(y))$ nie ma zer na brzegu.
Niezmienniczość stopnia daje
$\deg(\id-F_t,U,0)=\deg(V,U,0)$.
Jeśli $\D_pX$ jest odwracalna, lokalny stopień pola
wynosi znak jej wyznacznika. Dodatkowo
$\D_pf_t=I-t\D_pX$ (po identyfikacji przestrzeni stycznej),
więc znak $\det(I-\D_pf_t)$ jest ten sam.
\end{proof}

\begin{wniosek}[Dowód Poincarégo--Hopfa przez Lefschetza]
Niech $M$ będzie zamkniętą, zorientowaną, gładką
$n$-rozmaitością, $n\ge1$, a $X$ gładkim polem z izolowanymi
zerami. Wtedy $\sum_{X(p)=0}\operatorname{ind}_p(X)=\chi(M)$.
\end{wniosek}
\begin{proof}
Dla małego $t$ odwzorowanie $f_t=\exp(-tX)$ jest globalnie
określone i homotopijne z identycznością poprzez
$s\mapsto\exp(-stX)$. Jego punkty stałe są dokładnie zerami $X$.
Dla zer nieosobliwych Twierdzenie~\ref{thm:lefschetz-hopf-v10}
oraz Lemat~\ref{lem:field-fixed-v11} dają
$\sum_p\operatorname{ind}_p(X)=L(f_t)=\chi(M)$.
W przypadku izolowanych zer osobliwych stosujemy małe,
lokalne perturbacje pola, utrzymując je bez zmian na brzegach
rozłącznych kul otaczających zera. Transwersalność do
przekroju zerowego zapewnia nieosobliwość nowych zer,
a homotopijna niezmienniczość stopnia gwarantuje, że suma
ich indeksów w każdej kuli równa się staremu indeksowi.
\end{proof}

\section{Wersja z brzegiem: pole skierowane na zewnątrz}
\begin{definicja}[Pole skierowane ściśle na zewnątrz]
Niech $M$ będzie zwartą gładką rozmaitością z brzegiem,
a $\nu$ zewnętrznym jednostkowym normalnym polem
wzdłuż $\partial M$ dla ustalonej metryki.
Pole styczne $X$ na $M$ jest \emph{ściśle skierowane na zewnątrz},
jeśli $\langle X(x),\nu(x)\rangle>0$ dla każdego
$x\in\partial M$. Warunek nie zależy od wybranej metryki:
wyraża dodatniość składowej normalnej w koorientacji brzegu.
\end{definicja}

\begin{twierdzenie}[Poincaré--Hopf dla rozmaitości z brzegiem]
\label{thm:PH-boundary-v11}
Niech $M$ będzie zwartą zorientowaną gładką
$n$-rozmaitością z brzegiem, $n\ge1$.
Jeśli gładkie pole $X\in\mathfrak X(M)$ jest ściśle skierowane
na zewnątrz na $\partial M$ i ma izolowane zera we wnętrzu,
to zbiór tych zer jest skończony oraz
\[
\sum_{p\in\operatorname{int}M:\,X(p)=0}\operatorname{ind}_p(X)=\chi(M).
\]
\end{twierdzenie}
\begin{proof}
Zera nie mogą zbliżać się do brzegu, bo pole ma tam
niezerową składową normalną. Zatem tworzą zwarty dyskretny
podzbiór wnętrza i jest ich skończenie wiele.
Najpierw zakładamy, że są nieosobliwe.
Wybierzmy metrykę i pole przedłużone do otoczenia
$M$ w jej podwojeniu (lub innej rozmaitości bez brzegu).
Tam definiujemy mapę wykładniczą. Przedłużenie nie
musi zachowywać warunku skierowania na zewnątrz
poza $M$, ponieważ rozważamy punkty $x\in M$.
Dla dostatecznie małego $t>0$ odwzorowanie
$f_t(x)=\exp_x(-tX(x))$ jest określone na całym $M$
i przyjmuje wartości w $M$: wewnątrz wynika to z małości
$t$, a przy brzegu wektor $-X$ jest ściśle skierowany
do wnętrza; przez zwartość warunek jest jednostajny
na kołnierzu. Mamy $f_t\simeq\id_M$ przez tę samą
rodzinę z parametrem $st$, a punkty stałe $f_t$ są
wyłącznie wewnętrznymi zerami $X$.

Dla ścisłego zastosowania twierdzenia indeksowego
przechodzimy do podwojenia $DM$ rozmaitości $M$.
Niech $j:M\hookrightarrow DM$ oznacza pierwszą kopię,
a $r:DM\to M$ odwzorowanie składające obie kopie
na pierwszą. Spełnia ono $r\circ j=\id_M$.
Odwzorowanie $F=j\circ f_t\circ r:DM\to DM$ ma dokładnie
te same punkty stałe co $f_t$; wszystkie leżą we wnętrzu
pierwszej kopii. Na ich otoczeniach $F$ jest gładkie.
Po wygładzeniu $F$ poza tymi otoczeniami, względnie do nich,
możemy zastosować formułę Lefschetza--Hopfa na zamkniętej
rozmaitości $DM$. Wygładzenie nie tworzy nowych punktów
stałych, gdyż na zwartym dopełnieniu małych otoczeń
odległość punktu od jego obrazu ma dodatnie minimum.

Cykliczność śladu dla $j_*,r_*$ i $r_*j_*=\id$ daje
$L(F)=L(f_t)=L(\id_M)=\chi(M)$.
Lokalne indeksy $F$ i $f_t$ są identyczne, zaś
Lemat~\ref{lem:field-fixed-v11} identyfikuje je z indeksami
pola $X$. Otrzymujemy żądaną równość.
Dla izolowanych zer osobliwych perturbujemy $X$ w
skończenie wielu rozłącznych kulach we wnętrzu,
zachowując warunek na brzegu i indeksy lokalne,
po czym stosujemy przypadek nieosobliwy.
\end{proof}

\begin{przyklad}[Dysk, kula i pole radialne]
Na $D^n\subset\R^n$ pole $X(x)=x$ jest ściśle zewnętrzne
na brzegu i ma jedno zero w $0$, którego indeks wynosi $+1$.
Stąd $\chi(D^n)=1$. Na kuli z usuniętym wnętrzem mniejszej
współśrodkowej kuli pole $X(x)=x$ nie jest zewnętrzne
na wewnętrznej składowej brzegu; nie można więc bez
zmiany hipotez stosować powyższego twierdzenia.
\end{przyklad}

\begin{przyklad}[Pole gradientowe funkcji Morse'a]
Niech $M$ będzie zwartą rozmaitością bez brzegu,
a $f:M\to\R$ funkcją Morse'a. Dla $X=\operatorname{grad}f$
każde zero $p$ jest nieosobliwe i
$\operatorname{ind}_p(X)=(-1)^{\lambda(p)}$.
Dlatego twierdzenie Poincarégo--Hopfa daje
$\chi(M)=\sum_p(-1)^{\lambda(p)}$.
Jest to ta sama równość, którą otrzymaliśmy algebraicznie
z kompleksu Morse'a, ale interpretacja lokalna jest inna.
\end{przyklad}

\section{Co się zmienia, gdy pole jest styczne do brzegu?}
\begin{uwaga}
Jeśli $X$ jest tylko styczne do $\partial M$ albo jego
składowa normalna zmienia znak, suma indeksów zer wewnętrznych
nie musi równać się $\chi(M)$. Trzeba uwzględnić warunek
brzegowy, np. stopień części normalizowanej pola na brzegu
lub odpowiedni indeks brzegowy. Nie podajemy niepoprawnego
wzoru bez ustalenia konwencji dla tych indeksów.
\end{uwaga}


% Iteracja XI: samodzielne wprowadzenie do nośników zwartych i Borel--Moore'a.


% Rozdział do wstawienia do głównego skryptu; korzysta z jego preambuły.
